%--------------------------------------------------------------------------------
%
%
%--------------------------------------------------------------------------------
\pagebreak
\section*{PRB — Probe Virtual Address}
\addcontentsline{toc}{section}{PRB — Probe Virtual Address}

\begin{iPageDescEntry}{Syntax}
    \texttt{PRB.A RegR, RegB, RegA} \\
    \texttt{PRB .R/W/X RegR, RegB}
\end{iPageDescEntry}

\begin{iPageDescEntry}{Format}
    The \texttt{PRB} instruction uses both the following instruction format: \\

    \begin{flushleft}
    \drawInstrFormat
        {0,1,3,5,6.5,8.5,9.5,11.5,16}
        {\OpCodeGrpSYS, \OpCodePRB, RegR, 0, RegB, mode, RegA, 0}
    \end{flushleft}
   
\end{iPageDescEntry}

\begin{iPageDescEntry}{Description}
    The \texttt{PRB} instruction probes the Translation Lookaside Buffer (TLB) 
    to determine whether the effective address in \texttt{RegB} has a valid mapping. 
    The result of the probe is returned in \texttt{RegR}. A value of \texttt{0} indicates that no mapping exists, while a \texttt{1} indicates a valid mapping.  
    
    The test option is encoded in the mode field. An instruction option of \texttt{R} encodes as zero, a \texttt{W} encodes as one and \texttt{X} encodes as two. A mode field value of three indicates that the instruction option is encoded in right most two bits of \texttt{RegA}. Here a value of 3 is undefined.
\end{iPageDescEntry}

\begin{iPageDescOpEntry}{Operation}
\begin{lstlisting}[style=iPageOpStyle]
RegR <- probeVirtualAddress( RegB, RegA.[1..0]);    // Reg mode
RegR <- probeVirtualAddress( RegB, instr.[mode] );  // Imm mode
\end{lstlisting}
\end{iPageDescOpEntry}

\begin{iPageDescEntry}{Exceptions}
    - Non-access TLB Trap \\
    - Privileged Operation Trap
\end{iPageDescEntry}

\begin{iPageDescEntry}{Notes}
    Typically used by the operating system for virtual memory management.
\end{iPageDescEntry}
